\documentclass[12pt]{article}
\usepackage[T1]{fontenc}
\usepackage[utf8]{inputenc}
\usepackage{lmodern}
\usepackage{textcomp}
\usepackage{enumitem}
\usepackage{geometry}
\geometry{margin=1in}
\begin{document}
\begin{center}
Answer Key - Health Sciences - Graduate Level Assessment 17 \
\small (Answers are ordered exactly as in the question paper.)
\end{center}

\section*{Multiple Choice}
\begin{enumerate}[label=\arabic*.]
\item \textbf{Answer:} A
\\ \textit{Options:} A) If the patient chooses it.; B) If the patient is bedridden.; C) If urine does not need to be measured.; D) If the patient has bladder sensation.; E) If the patient suffers from urinary incontinence.
\\ \textit{Note:} Correct option: A - If the patient chooses it.
\item \textbf{Answer:} A
\\ \textit{Options:} A) loss of taste over the anterior two-thirds of the tongue.; B) complete loss of taste sensation.; C) paralysis of the muscles of the throat.; D) paralysis of the muscles of the jaw.; E) loss of somaesthetic sensation over the posterior two thirds of the tongue.
\\ \textit{Note:} Correct option: A - loss of taste over the anterior two-thirds of the tongue.
\item \textbf{Answer:} A
\\ \textit{Options:} A) Zellweger syndrome; B) Maple syrup urine disease; C) Hemophilia; D) Medium chain acyl-CoA dehydrogenase deficiency; E) Huntington's disease
\\ \textit{Note:} Correct option: A - Zellweger syndrome
\item \textbf{Answer:} A
\\ \textit{Options:} A) motor and autonomic neuronal processes.; B) both autonomic and motor neuronal processes.; C) both sensory and motor neuronal processes.; D) sensory and autonomic neuronal processes.; E) motor neuronal processes.
\\ \textit{Note:} Correct option: A - motor and autonomic neuronal processes.
\item \textbf{Answer:} E
\\ \textit{Options:} A) diaphragm.; B) diaphragm and pericardium.; C) diaphragm, parietal pleura, pericardium and abdominal muscles.; D) diaphragm, parietal pleura, pericardium and lungs.; E) diaphragm, parietal pleura and pericardium.
\\ \textit{Note:} Correct option: E - diaphragm, parietal pleura and pericardium.
\end{enumerate}

\section*{True / False}
\begin{enumerate}[label=\arabic*.]
\setcounter{enumi}{5}
\item \textbf{Answer:} False
\\ \textit{Options:} A) True; B) False
\\ \textit{Note:} LLM-generated false statement. Correct answer: 'Are more valid'.
\item \textbf{Answer:} False
\\ \textit{Options:} A) True; B) False
\\ \textit{Note:} LLM-generated false statement. Correct answer: 'can be carried out by the patient.'.
\item \textbf{Answer:} True
\\ \textit{Options:} A) True; B) False
\\ \textit{Note:} LLM-generated true statement based on MCQ.
\item \textbf{Answer:} True
\\ \textit{Options:} A) True; B) False
\\ \textit{Note:} LLM-generated true statement based on MCQ.
\item \textbf{Answer:} True
\\ \textit{Options:} A) True; B) False
\\ \textit{Note:} LLM-generated true statement based on MCQ.
\end{enumerate}

\section*{Long Form Response}
\begin{enumerate}[label=\arabic*.]
\setcounter{enumi}{10}
\item \textbf{Answer:} Autosomal recessive. Explain why this is the correct answer and provide supporting evidence. Reference key facts or concepts that support this choice.
\\ \textit{Note:} Long-form derived from MCQ; award credit for stating the correct idea and giving supporting reasoning.
\item \textbf{Answer:} Thalassemia. Explain why this is the correct answer and provide supporting evidence. Reference key facts or concepts that support this choice.
\\ \textit{Note:} Long-form derived from MCQ; award credit for stating the correct idea and giving supporting reasoning.
\end{enumerate}

\vfill
\noindent\textit{End of Answer Key}
\end{document}